\documentclass{article}
\usepackage{mate}

\title{Algebra liniara}
\author{}
\date{}

\begin{document}

\maketitle

\newcommand{\colvec}[1]{\begin{bmatrix}#1\end{bmatrix}}

\section{Spatii vectoriale}
\begin{enumerate}

  \definitie Se numeste $\mathbb{K}\term{-spatiu vectorial}$ structura algebrica ${(V, \mathbb{K}, +, \cdot)}$, unde $(\mathbb{K}, +, \cdot)$ este un corp comutativ, $V \neq \varnothing$ si $(V, +)$ e grup abelian, care verifica axiomele:
  \begin{enumerate}
    \item $\alpha (x + y) = \alpha x + \alpha y, \forall \alpha \in \mathbb{K}, x, y \in V$. 
    \item $(\alpha + \beta)x = \alpha x + \beta x, \forall \alpha, \beta \in \mathbb{K}, x \in V$
    \item $\alpha \cdot (\beta x) = (\alpha \cdot \beta) x, \forall \alpha, \beta \in \mathbb{K}, x \in V$.
    \item $1_{\mathbb{K}} \cdot x = x, \forall x \in V$.

  \end{enumerate}
  Elementele din $V$ se numesc \term{vectori}, iar elementele din $\mathbb{K}$ se numesc \term{scalari}.
  
  \exemplu Orice corp comutativ $\mathbb{K}$ este si $\mathbb{K}$-spatiu vectorial.
  \exemplu Daca $\mathbb{K}$ este corp comutativ si $n \in \mathbb{N}^*$, atunci $\mathbb{K}^n = \left\{ \colvec{x_1 \\ x_2 \\ \vdots \\ x_n} \middle\vert x_1, x_2, \dots, x_n \in \mathbb{K} \right\}$ este $\mathbb{K}$-spatiu vectorial, cu $\colvec{x_1 \\ x_2 \\ \vdots \\ x_n} + \colvec{y_1 \\ y_2 \\ \vdots \\ y_n} = \colvec{x_1 + y_1 \\ x_2 + y_2 \\ \vdots \\ x_n + y_n}$ si $\alpha \colvec{x_1 \\ x_2 \\ \vdots \\ x_n} = \colvec{\alpha x_1 \\ \alpha x_2 \\ \vdots \\ \alpha x_n}$.
  
  \exemplu $\mathbb{C}$ este spatiu vectorial real, dar si complex.
  \exemplu $\mathcal{M}_{m, n}(\mathbb{K})$ este $\mathbb{K}$-spatiu vectorial.
  \exemplu $\mathbb{K}_n [X] = \left\{ P \in \mathbb{K}[X] \middle\vert grad P \leq n \right\}$ este $\mathbb{K}$-spatiu vectorial.
  
  \exemplu $\widetilde{\mathbb{K}}_n [X] = \left\{ P \in \mathbb{K}[X] \middle\vert grad P = n \right\}$ NU este $\mathbb{K}$-spatiu vectorial, deoarece $grad(P + (-P)) = -\infty$, desi $P$ si $-P$ $\in \widetilde{\mathbb{K}}_n[X]$.

  \proprietate Fie $(V, \mathbb{K}, +, \cdot)$ un $\mathbb{K}$-spatiu vectorial. Atunci:
  \begin{enumerate}
    \item $0_{\mathbb{K}} \cdot x = 0_V$.
    \item $\alpha \cdot 0_V = 0_V$.
    \item $\alpha \cdot x = 0_V \Leftrightarrow \alpha = 0_{\mathbb{K}} \text{ sau } x = 0_V$.
    \item $(-1_{\mathbb{K}}) \cdot x = -x$
  \end{enumerate}


  \begin{tcolorbox}
  \definitie Fie $(V, \mathbb{K}, +, \cdot)$ un $\mathbb{K}$-spatiu vectorial. 
  \begin{enumerate}
    \item Se numeste \term{combinatie liniara} a vectorilor \\ $v_1, v_2, \dots, v_n \in V$ expresia $\alpha_1 v_1 + \alpha_2 v_2 + \dots + \alpha_n v_n$, unde $\alpha_1, \alpha_2, \dots, \alpha_n \in \mathbb{K}$.

  \item Se numeste \term{sistem de generatori} (sau SG) pentru $V$ o submultime $S \subseteq V$ cu proprietatea ca $\forall v \in V, \\ \exists v_1, v_2, \dots v_n \in S$ si $\alpha_1, \alpha_2, \dots, \alpha_n \in \mathbb{K}$ a.i $v = \alpha_1 v_1 + \alpha_2 v_2 + \dots + \alpha_n v_n$ (orice vector din $V$ este o combinatie liniara de vectori din $S$).

  \item Se numeste \term{sistem liniar independent} (sau SLI) o submultime $S = \{v_1, v_2, \dots, v_n\} \subseteq V$ cu proprietatea ca $\alpha_1 v_1 + \alpha_2 v_2 + \dots + \alpha_n v_n = 0_V \Rightarrow \alpha_1 = \alpha_2 = \dots = \alpha_n = 0_{\mathbb{K}}$ (niciun vector din $S$ nu poate fi combinatie liniara nenula de alti vectori din $S$).
  \item Se numeste \term{sistem liniar dependent} (sau SLD) o submultime $S \subseteq V$ care nu este sistem liniar independent (exista un vector in $S$ care poate fi combinatie liniara nenula de alti vectori din $S$).


\end{enumerate}
\end{tcolorbox}

\observatie Daca $S$ este SLI, atunci $0_V \notin S$.

\propozitie Fie $(V, \mathbb{K}, +, \cdot)$ un $\mathbb{K}$-spatiu vectorial si $S \subseteq V$ un SLI. Atunci orice submultime a lui $S$ este SLI.

\propozitie Fie $(V, \mathbb{K}, +, \cdot)$ un $\mathbb{K}$-spatiu vectorial si $S = \{v_1, v_2, \dots, v_n\}$ un SLD de vectori nenuli. Atunci $\exists j \in \{2, 3, \dots, n\}$ si $a_1, a_2, \dots, a_{j-1} \in \mathbb{K}$ a.i $v_j = \sum_{k=1}^{j-1} a_k v_k$. (Cel putin un vector poate fi combinatie liniara nenula a vectorilor dinaintea sa)

\begin{tcolorbox}
  \definitie Fie $(V, \mathbb{K}, +, \cdot)$ un $\mathbb{K}$-spatiu vectorial. Se numeste \term{baza} a lui $V$ o submultime $B \subseteq V$ care este SLI si SG (sistem de generatori) pentru $V$.
  \teorema (Existenta bazei) Orice spatiu vectorial nenul finit generat are cel putin o baza.
\end{tcolorbox}
  \observatie Un spatiu vectorial este finit generat daca exista o multime finita care este SG pentru acel spatiu.
  \exemplu Baza canonica a lui $\mathbb{R}^n = \mathcal{M}_{n, 1}(\mathbb{R})$ este formata din coloanele lui $I_n$.
  \[
    (Bc)_{\mathbb{R}^n} = \left\{
      c_1 = \colvec{1 \\ 0 \\ \vdots \\ 0},
      \dots,
      c_n = \colvec{0 \\ 0 \\ \vdots \\ 1}
    \right\}.
  \]

  \exemplu Baza canonica a lui $\mathcal{M}_{m, n}(\mathbb{R})$ este $\{E_{ij} \mid 1 \leq i \leq m,\ 1 \leq j \leq n\}$, unde $E_{ij} = [a_{k\ell}]_{m \times n}$ si
  \[
    a_{k\ell} =
    \begin{cases}
      1, & \text{daca } k = i \text{ si } \ell = j, \\
      0, & \text{altfel}.
    \end{cases}
  \]
  In particular,
  \[
    (Bc)_{\mathcal{M}_{2, 3}(\mathbb{R})} =
    \left\{
      \colvec{1 & 0 & 0 \\ 0 & 0 & 0},
      \colvec{0 & 1 & 0 \\ 0 & 0 & 0},
      \colvec{0 & 0 & 1 \\ 0 & 0 & 0},
      \colvec{0 & 0 & 0 \\ 1 & 0 & 0},
      \colvec{0 & 0 & 0 \\ 0 & 1 & 0},
      \colvec{0 & 0 & 0 \\ 0 & 0 & 1}
    \right\}.
  \]
  \exemplu $\{1, X, X^2, \dots, X^n\}$ este baza canonica pentru $\mathbb{R}_n[X]$.
  \exemplu $\{1, i\}$ este o baza pentru $\mathbb{C}$ privit ca spatiu vectorial real.
  \exemplu Daca $z_0 \in \mathbb{C}^*$, atunci $\{z_0\}$ este o baza pentru $\mathbb{C}$ privit ca spatiu vectorial complex.

  \definitie Fie $(V, \mathbb{K}, +, \cdot)$ un $\mathbb{K}$-spatiu vectorial si $B = \{v_1, v_2, \dots, v_n\}$ o baza a lui $V$. Atunci $\forall v \in V$, $\exists \alpha_1, \alpha_2, \dots, \alpha_n \in \mathbb{K}$ a.i $v = \alpha_1 v_1 + \alpha_2 v_2 + \dots + \alpha_n v_n$. Scalarii $\alpha_1, \alpha_2, \dots, \alpha_n$ se numesc \term{coordonatele vectorului $v$ in baza $B$}, in ordinea $v_1, v_2, \dots, v_n$.

  \notatie $[v]_B = \colvec{\alpha_1 \\ \alpha_2 \\ \vdots \\ \alpha_n}$.

  \propozitie Fie $(V, \mathbb{K}, +, \cdot)$ un $\mathbb{K}$-spatiu vectorial finit generat si $B$ o baza a lui $V$. Atunci coordonatele oricarui vector $v \in V$ in baza $B$ sunt unice.

  \lema (Schimbului) Fie $(V, \mathbb{K}, +, \cdot)$ un $\mathbb{K}$-spatiu vectorial nenul, $S = \{v_1, v_2, \dots, v_m\} \subseteq V \setminus \{0_V\}$ un SG pentru $V$ si $\{w_1, w_2, \dots, w_n\} \subseteq V$ un SLI. Atunci:
  \begin{enumerate}
    \item $n \leq m$.
    \item Dupa o eventuala renumerotare a vectorilor din $S$, $\{w_1, w_2, \dots, w_n, v_{n+1}, \dots, v_m\}$ este un SG pentru $V$.
  \end{enumerate}

  \begin{tcolorbox}
    \teorema (Cardinalul bazei) Intr-un spatiu vectorial nenul finit generat, toate bazele sunt finite si au acelasi numar de vectori. Acest numar se numeste \term{dimensiunea} spatiului vectorial si se noteaza $\dim_{\mathbb{K}} V$.
  \end{tcolorbox}

  \observatie $\dim_{\mathbb{R}} \mathbb{C} = 2$, iar $\dim_{\mathbb{C}} \mathbb{C} = 1$.

\end{enumerate}

\end{document}
